\section{Обзор существующих подходов и решений}

\subsection{Классические модели управления потоками}

В процессе развития параллельного программирования сложилось несколько
устойчивых моделей управления потоками, каждая из которых представляет
собой определённый компромисс между простотой реализации, накладными
расходами и эффективностью использования ресурсов.

Исторически первой и наиболее простой является модель «поток на
задачу» (thread-per-task), при которой для каждой входящей задачи
создаётся отдельный поток, уничтожаемый по завершении работы.
Достоинством данной модели является концептуальная простота:
разработчику не требуется явно управлять жизненным циклом потоков
или координировать распределение задач. Недостатки этой модели ---
высокие накладные расходы на создание и уничтожение потоков.

Следующим шагом в развитии стала модель пула потоков с единой
глобальной очередью (single global queue thread pool). В данной модели
при запуске приложения создаётся фиксированное число потоков,
равное, как правило, числу логических ядер процессора. Все потоки
разделяют одну очередь задач: производитель помещает задачи в конец
очереди, а потоки-потребители извлекают их из начала. Данная модель
устраняет накладные расходы на создание и уничтожение потоков,
однако порождает новую проблему: разделяемая очередь становится
точкой конкуренции. При большом числе потоков операции постановки
и извлечения задач требуют синхронизации посредством мьютекса или
атомарных операций, что ограничивает масштабируемость
системы~\cite{Williams2019}.

Для устранения конкуренции за единую очередь была предложена модель
пула потоков с локальными очередями (per-thread queue thread pool).
В данной архитектуре каждый поток владеет собственной очередью задач.
Новые задачи назначаются потокам по некоторому правилу --- например,
по принципу Round Robin или путём хеширования идентификатора задачи.
Данный подход снижает конкуренцию, однако приводит к статическому
распределению задач, которое не адаптируется к реальной загрузке
потоков. В результате возможна ситуация, когда одни потоки перегружены,
тогда как другие простаивают~\cite{Herlihy2020}.

Попыткой совместить преимущества обеих моделей стала архитектура
с иерархическими очередями (hierarchical queue). В ней глобальная
очередь выступает резервуаром задач, а каждый поток дополнительно
располагает локальным буфером для хранения небольшого числа задач.
Поток в первую очередь берёт задачи из локального буфера, обращаясь
к глобальной очереди лишь при его опустошении. Данный подход снижает
частоту синхронизации, но не решает проблему неравномерной загрузки
потоков~\cite{Williams2019}.

Отдельного рассмотрения заслуживает модель реакторов (reactor pattern)
и её расширение --- модель SEDA (Staged Event-Driven Architecture).
В модели реактора единственный поток обрабатывает события ввода-вывода
и диспетчеризирует задачи, что позволяет избежать накладных расходов
на синхронизацию, однако не позволяет эффективно задействовать
несколько ядер. Модель SEDA разделяет обработку на несколько
последовательных стадий, каждая из которых обслуживается собственным
пулом потоков, что обеспечивает лучшую масштабируемость и
управляемость нагрузки~\cite{Welsh2001}.

Анализ перечисленных моделей позволяет выявить фундаментальное
противоречие, характерное для классических подходов: снижение
конкуренции за разделяемые структуры данных достигается ценой
ухудшения балансировки нагрузки, и наоборот. Данное противоречие
является основной мотивацией для разработки механизма work-stealing,
который стремится одновременно минимизировать синхронизацию и
обеспечить динамическую балансировку нагрузки.

\subsection{Механизм work-stealing: история возникновения
и теоретические основы}

Механизм work-stealing (кража задач) представляет собой алгоритм
динамического распределения задач между потоками, при котором
незагруженный поток вправе забирать («красть») задачи из очередей
других, более загруженных потоков. Данный подход позволяет совместить
низкую конкуренцию за локальные очереди с автоматической балансировкой
нагрузки, не требующей централизованного координатора.

Истоки механизма work-stealing восходят к исследованиям в области
параллельных вычислений, проводившимся в 1980-е годы. Одной из первых
систем, реализовавших схожие принципы, стал язык Multilisp,
разработанный Робертом Хэлстедом в Массачусетском технологическом
институте в 1985 году~\cite{Halstead1985}. В Multilisp параллельные
вычисления организовывались посредством конструкции \texttt{pcall},
а распределение задач между процессорами осуществлялось динамически.

Однако подлинной теоретической основой современного work-stealing
стала система Cilk, разработанная в MIT в 1990-е годы группой под
руководством Чарльза Лейзерсона. В основополагающей работе Блюмофе
и Лейзерсона~\cite{Blumofe1999} была формализована модель
многопоточных вычислений и доказана оптимальность жадного алгоритма
work-stealing с точки зрения теоретической эффективности.

В модели Блюмофе и Лейзерсона многопоточное вычисление представляется
в виде ориентированного ациклического графа (DAG), вершины которого
соответствуют элементарным инструкциям, а рёбра задают отношения
зависимости. Для данной модели вводятся два ключевых параметра:
работа $T_1$ --- суммарное число инструкций при выполнении
на одном процессоре, и \textit{глубина} (критический путь) $T_\infty$
--- длина наидлиннейшей цепочки зависимых инструкций.

Авторы доказали, что при выполнении на $P$ процессорах алгоритм
work-stealing обеспечивает время выполнения:

\begin{equation}
    T_P \leq \frac{T_1}{P} + O(T_\infty)
    \label{eq:work-stealing-bound}
\end{equation}

Данная оценка является оптимальной с точностью до константного
множителя: первое слагаемое отражает теоретический минимум,
обусловленный объёмом работы, второе --- неустранимые потери
из-за последовательных зависимостей~\cite{Blumofe1999}.

Помимо оценки времени выполнения, Блюмофе и Лейзерсон доказали,
что алгоритм work-stealing является оптимальным по использованию
памяти: потребление памяти при выполнении на $P$ процессорах
не превышает $P$-кратного потребления при последовательном выполнении.
Данное свойство принципиально важно для практических приложений,
поскольку бесконтрольный рост потребления памяти является типичной
проблемой наивных реализаций параллельных вычислений.

Ключевым структурным элементом алгоритма work-stealing является
двусторонняя очередь (double-ended queue, deque), которой владеет
каждый поток. Работа с данной структурой данных подчиняется следующим
правилам:

\begin{itemize}
    \item поток-владелец помещает новые задачи в нижний конец очереди
    (push) и извлекает задачи для выполнения также из нижнего конца
    (pop), то есть работает по принципу стека (LIFO);
    \item поток-«вор», не имеющий собственных задач, извлекает задачи
    из верхнего конца очереди жертвы (steal), то есть работает
    по принципу очереди (FIFO).
\end{itemize}

Такое асимметричное использование deque преследует две цели.
Во-первых, доступ владельца к нижнему концу в большинстве случаев
не требует синхронизации, поскольку конкуренция возможна лишь
в граничной ситуации, когда очередь содержит ровно одну задачу.
Во-вторых, воры забирают задачи с верхнего конца, то есть наиболее
«старые» задачи, порождённые ранее всего. В рекурсивных алгоритмах
типа «разделяй и властвуй» такие задачи, как правило, являются
наиболее крупными и предоставляют больше всего работы, что делает
кражу максимально эффективной~\cite{Acar2000}.

Важным практическим достижением стала работа Чейза и Лева~\cite{Chase2005},
в которой была предложена реализация deque для work-stealing на основе
динамически расширяемого кольцевого буфера. Данная реализация требует
синхронизации (посредством инструкции CAS) лишь в случае конкуренции
между владельцем и вором за последнюю задачу в очереди, обеспечивая
тем самым практически неблокирующее поведение в типичных сценариях.

После публикации теоретических работ механизм work-stealing был
реализован в ряде промышленных систем и языков программирования.
Библиотека Intel Threading Building Blocks (TBB), выпущенная в 2006
году, стала одной из первых широко используемых реализаций
work-stealing для языка C++~\cite{Blumofe1999}. В языке Java механизм
work-stealing был включён в стандартную библиотеку в виде класса
\texttt{ForkJoinPool}, появившегося в версии Java~7 в 2011 году.
В языке C++ стандарт не предоставляет встроенной реализации
work-stealing, что обусловливает актуальность разработки собственных
эффективных реализаций~\cite{Williams2019}.

Таким образом, механизм work-stealing сочетает в себе строгое
теоретическое обоснование и практическую эффективность, что делает
его одним из наиболее перспективных подходов к управлению задачами
в параллельных системах.